\subsection{Onto clocks and the exact derivative estimates}
\label{bcs:clocks-section}
Slow forcing must still complete every finite logical step.  The relevant
condition is that the new clock maps the nonnegative half-line onto itself.
The following elementary calculation applies to our transverse pulse
templates, including their finite loading intervals.

\begin{lemma}[Onto clocks]\label{bcs:clock}
Let $W(s,x)$ be a smooth divergence-free field on a flat torus, with
zero mean, $W(0)=0$, bounded mixed derivatives and
$(W\cdot\nabla)W=0$.  Let $g:[0,\infty)\to[0,\infty)$ be
smooth, nondecreasing and onto, with $g(0)=0$.  Then
\begin{align*}
 u(t,x)&=g'(t)W(g(t),x),\\
 f(t,x)&=g''(t)W(g(t),x)+(g'(t))^2W_s(g(t),x)
                         -\nu g'(t)\Delta W(g(t),x)
\end{align*}
give a mean-zero solenoidal force and a solution with zero pressure and
zero initial velocity.  Its flow is $X_u(t,a)=X_W(g(t),a)$, so every
all-time reachability event is unchanged.  For $g(t)=\log(1+t)$ and
every $k,\alpha$,
\[
 \|\partial_t^kD_x^\alpha u(t)\|_\infty+
 \|\partial_t^kD_x^\alpha f(t)\|_\infty
 \le C_{k,\alpha}(1+t)^{-1-k}.
\]
All these derivatives are square integrable in time in supremum norm,
and in space--time.
\end{lemma}
\begin{proof}
The chain rule gives $f=u_t-\nu\Delta u$.  Divergence, mean zero,
and vanishing self-advection persist under multiplication by the scalar
$g'(t)$ and time composition.  The same chain rule proves the trajectory
identity by ODE uniqueness.  Onto ensures that every finite logical time
has a finite physical preimage, proving both implications for the event.
The energy comparison in Lemma~\ref{sv:force} applies to this explicit
solution on every finite interval, without a periodicity assumption.

For the logarithmic choice put $a=(1+t)^{-1}$ and $s=\log(1+t)$.
The velocity is $aW(s)$ and the force is
$-\nu a\Delta W(s)+a^2(W_s-W)(s)$.
For each positive integer $r$,
\[
 \partial_t[a^rA(s,x)]=a^{r+1}(\partial_s-r)A(s,x).
\]
Repeated differentiation and the bounded template derivatives give the
stated estimate.  Its square is integrable, and the torus has finite
volume.  These are estimates for the precomputed force formula and
require no simulation of the instructions.
\end{proof}

\begin{proposition}[Cube-root and fourth-root clocks]
\label{bcs:power-clock}
Under the hypotheses on $W$ in Lemma~\ref{bcs:clock}, choose
$g(t)=(1+t)^q-1$, with $q=1/3$ or $1/4$, and put $\beta=1-q$.
The resulting force and velocity have bounded mixed derivatives and
\[
 \|\partial_t^kD_x^\alpha u(t)\|_\infty+
 \|\partial_t^kD_x^\alpha f(t)\|_\infty
 \le C_{k,\alpha}(1+t)^{-\beta}
\]
for every fixed $k,\alpha$.  Each such derivative is square integrable
in time in spatial supremum norm.
\end{proposition}
\begin{proof}
Every positive derivative $g^{(j)}$ is bounded and is
$O_j((1+t)^{-\beta})$.  Every differentiated term of $g'W(g)$
contains at least one such factor and a bounded derivative of $W$.
One factor supplies the decay and the remaining factors are bounded.
The identity $f=u_t-\nu\Delta u$ proves the same estimate for the
force.  Finally $2\beta>1$, and each clock is increasing and onto.
\end{proof}

The force obtained after slowing is a different prescription from the
unslowed periodic force.  In particular these conclusions do not combine
decay with physical period one.  They also do not imply decay faster than
every power of time for an arbitrary repeated processor: a speed bounded
by $C(1+t)^{-2}$ has only finite accumulated logical time.
The same obstruction appears in the viscous example of
Cardona--Miranda--Peralta-Salas--Presas~\cite[Section~6B]{CMPP}:
their exponentially damped velocity traverses only a finite interval
of the underlying stationary trajectory.  The onto condition above
ensures that slowing preserves every finite stage of the computation.
